\documentclass[11pt]{article}
\usepackage[margin=1in]{geometry}
\usepackage{amsmath, amssymb, amsthm, enumitem}

\title{\textbf{GRE Mathematics Subject Test}\\[4pt] Abstract Algebra Practice Problems}
\author{}
\date{}

\newcounter{prob}
\newcommand{\problem}{\stepcounter{prob}\medskip\noindent\textbf{Problem \theprob.}\ }

\begin{document}
\maketitle
\thispagestyle{empty}

%======================================================================
\section*{A. Groups: Basics \& Subgroups}
%======================================================================

\problem
Let $G$ be a group of order 35. Which of the following must be true?
\begin{enumerate}[label=(\alph*)]
\item $G$ is abelian.
\item $G$ is cyclic.
\item $G$ has an element of order 5.
\item All of the above.
\end{enumerate}

\problem
How many elements of order 2 does the symmetric group $S_3$ have?

\problem
Let $G$ be a group and $a\in G$ with $\operatorname{ord}(a)=12$. What is $\operatorname{ord}(a^8)$?

\problem
Which of the following groups are abelian? Select \textbf{all} that apply.
\begin{enumerate}[label=(\alph*)]
\item $S_3$
\item $\mathbb{Z}/6\mathbb{Z}$
\item $D_3$ (dihedral group of order 6)
\item The Klein four-group $V_4$
\end{enumerate}

\problem
Let $G$ be a group and $H\subseteq G$ defined by $H = \{g\in G : g^2 = e\}$. Is $H$ necessarily a subgroup of $G$?

\problem
What is the center $Z(S_3)$ of the symmetric group $S_3$?

%======================================================================
\section*{B. Cyclic Groups, Cosets \& Lagrange}
%======================================================================

\problem
How many generators does the cyclic group $\mathbb{Z}/20\mathbb{Z}$ have?

\problem
True or False: If $|G|=12$, then $G$ must have a subgroup of order 6.

\problem
Let $G=\mathbb{Z}/12\mathbb{Z}$. List all subgroups of $G$.

\problem
In the group $\mathbb{Z}/30\mathbb{Z}$, what is the order of the element $\bar{12}$ (i.e., $12 \bmod 30$)?

\problem
Let $G$ be a group of order 77. How many elements of order 7 does $G$ have?

%======================================================================
\section*{C. Normal Subgroups, Quotients \& Homomorphisms}
%======================================================================

\problem
Let $\phi\colon \mathbb{Z}\to\mathbb{Z}/6\mathbb{Z}$ be the natural projection $\phi(n)=n\bmod 6$. What is $\ker\phi$?

\problem
Let $G = \mathbb{Z}/12\mathbb{Z}$ and $N = \langle 4\rangle = \{0, 4, 8\}$. What is $G/N$ isomorphic to?

\problem
Let $\phi\colon S_3\to\mathbb{Z}/2\mathbb{Z}$ be the sign homomorphism ($\phi(\sigma)=0$ if $\sigma$ even, $1$ if odd). What is $\ker\phi$?

\problem
How many group homomorphisms are there from $\mathbb{Z}/6\mathbb{Z}$ to $\mathbb{Z}/4\mathbb{Z}$?

\problem
Let $G$ be a group and $H\le G$ with $[G:H]=2$. Which of the following must be true?
\begin{enumerate}[label=(\alph*)]
\item $H$ is normal in $G$.
\item $G/H\cong\mathbb{Z}/2\mathbb{Z}$.
\item $H$ is abelian.
\item (a) and (b) only.
\end{enumerate}

%======================================================================
\section*{D. Permutation Groups \& Sylow Theorems}
%======================================================================

\problem
In $S_5$, let $\sigma = (1\;2\;3)(4\;5)$. What is $\operatorname{ord}(\sigma)$?

\problem
How many elements of $S_4$ have order 3?

\problem
Is the permutation $(1\;2\;3\;4)(5\;6\;7)\in S_7$ even or odd?

\problem
Let $|G|=45=3^2\cdot 5$. How many Sylow 5-subgroups does $G$ have?

\problem
Let $|G|=30=2\cdot 3\cdot 5$. Which of the following is guaranteed by the Sylow theorems?
\begin{enumerate}[label=(\alph*)]
\item $G$ has a normal Sylow 5-subgroup.
\item $G$ has a normal Sylow 3-subgroup.
\item $G$ is abelian.
\end{enumerate}

%======================================================================
\section*{E. Rings \& Ideals}
%======================================================================

\problem
In $\mathbb{Z}/12\mathbb{Z}$, find all zero divisors.

\problem
Which of the following are integral domains? Select \textbf{all} that apply.
\begin{enumerate}[label=(\alph*)]
\item $\mathbb{Z}$
\item $\mathbb{Z}/6\mathbb{Z}$
\item $\mathbb{Z}/7\mathbb{Z}$
\item $\mathbb{Z}[x]$
\end{enumerate}

\problem
Is the ideal $(x)$ in $\mathbb{Z}[x]$ maximal, prime, both, or neither?

\problem
In the ring $\mathbb{Z}$, describe the quotient ring $\mathbb{Z}/(5)$.

\problem
Let $R=\mathbb{Z}[x]$. Is $(2,x)=\{2f(x)+xg(x) : f,g\in\mathbb{Z}[x]\}$ a principal ideal?

%======================================================================
\section*{F. Fields, Polynomials \& Miscellaneous}
%======================================================================

\problem
What is the characteristic of the field $\mathbb{F}_9$ (the field with 9 elements)?

\problem
Is $x^3 + x + 1$ irreducible over $\mathbb{F}_2 = \mathbb{Z}/2\mathbb{Z}$?

\problem
Find the degree of the field extension $[\mathbb{Q}(\sqrt[3]{2}):\mathbb{Q}]$.

\problem
How many abelian groups of order 36 are there, up to isomorphism?


%======================================================================
\newpage
\section*{Answer Key}
%======================================================================

\begin{enumerate}

\item (d) All of the above. $|G|=35=5\cdot 7$; since $5\nmid (7-1)=6$, any group of order $pq$ with $p<q$ and $p\nmid q-1$ is cyclic. Cyclic $\Rightarrow$ abelian. By Cauchy's theorem, $G$ has an element of order 5.

\item 3. The elements of order 2 in $S_3$ are the three transpositions: $(1\;2)$, $(1\;3)$, $(2\;3)$.

\item $\operatorname{ord}(a^8)=\frac{12}{\gcd(8,12)}=\frac{12}{4}=3$.

\item (b) and (d). $\mathbb{Z}/6\mathbb{Z}$ is cyclic hence abelian; $V_4\cong\mathbb{Z}/2\mathbb{Z}\times\mathbb{Z}/2\mathbb{Z}$ is abelian. $S_3$ and $D_3$ are both non-abelian (in fact $D_3\cong S_3$).

\item No, not necessarily. In an abelian group, $H$ is always a subgroup (since $(ab)^2=a^2b^2=e$). But in a non-abelian group it can fail: in $S_3$, $(1\;2)^2=e$ and $(2\;3)^2=e$ but $(1\;2)(2\;3)=(1\;2\;3)$ has order 3, so the set is not closed.

\item $Z(S_3)=\{e\}$, the trivial group. In $S_3$, no non-identity element commutes with every other element. (For example, $(1\;2)(1\;2\;3)=(2\;3)\neq(1\;3)=(1\;2\;3)(1\;2)$.)

\item $\varphi(20)=20\cdot(1-\frac{1}{2})\cdot(1-\frac{1}{5})=20\cdot\frac{1}{2}\cdot\frac{4}{5}=8$ generators.

\item False. The converse of Lagrange's theorem does not hold in general. However, for $|G|=12$ the statement actually depends on the specific group. The alternating group $A_4$ has order 12 but has no subgroup of order 6. So the answer is False.

\item The subgroups of $\mathbb{Z}/12\mathbb{Z}$ correspond to divisors of 12: $\langle 0\rangle=\{0\}$, $\langle 6\rangle=\{0,6\}$, $\langle 4\rangle=\{0,4,8\}$, $\langle 3\rangle=\{0,3,6,9\}$, $\langle 2\rangle=\{0,2,4,6,8,10\}$, and $\mathbb{Z}/12\mathbb{Z}$ itself. That is 6 subgroups of orders 1, 2, 3, 4, 6, 12.

\item $\operatorname{ord}(\bar{12})=\frac{30}{\gcd(12,30)}=\frac{30}{6}=5$.

\item $|G|=77=7\cdot 11$. Since $7\nmid(11-1)=10$, $G$ is cyclic ($\cong\mathbb{Z}/77\mathbb{Z}$). In a cyclic group of order 77, elements of order 7 are exactly those of the form $11k$ with $\gcd(k,7)=1$ for $k=1,\ldots,6$, giving $\varphi(7)=6$ elements of order 7.

\item $\ker\phi = 6\mathbb{Z} = \{\ldots, -12, -6, 0, 6, 12, \ldots\}$.

\item $|G/N|=12/3=4$. Since $G$ is cyclic and $G/N$ has order 4, $G/N$ is cyclic (quotient of a cyclic group is cyclic). So $G/N\cong\mathbb{Z}/4\mathbb{Z}$.

\item $\ker\phi = A_3 = \{e, (1\;2\;3), (1\;3\;2)\}$, the alternating group on 3 elements, which is cyclic of order 3.

\item The number of homomorphisms from $\mathbb{Z}/6\mathbb{Z}$ to $\mathbb{Z}/4\mathbb{Z}$ equals the number of elements $g\in\mathbb{Z}/4\mathbb{Z}$ with $6g=0$, i.e., $\operatorname{ord}(g)\mid 6$. Since $\operatorname{ord}(g)$ divides $|\mathbb{Z}/4\mathbb{Z}|=4$, we need $\operatorname{ord}(g)\mid\gcd(6,4)=2$. Elements of order dividing 2 in $\mathbb{Z}/4\mathbb{Z}$: $\bar{0}$ (order 1) and $\bar{2}$ (order 2). So there are \textbf{2} homomorphisms.

\item (d). Any subgroup of index 2 is normal (left and right cosets coincide), and $G/H$ has 2 elements so is isomorphic to $\mathbb{Z}/2\mathbb{Z}$. $H$ need not be abelian.

\item $\operatorname{ord}(\sigma)=\operatorname{lcm}(3,2)=6$.

\item The elements of order 3 in $S_4$ are exactly the 3-cycles. There are $\binom{4}{3}\cdot 2 = 8$ three-cycles in $S_4$.

\item $(1\;2\;3\;4)$ requires 3 transpositions (odd), $(5\;6\;7)$ requires 2 transpositions (even). Total: $3+2=5$ transpositions, so the permutation is \textbf{odd}.

\item $n_5\mid 9$ and $n_5\equiv 1\pmod{5}$. Divisors of 9: 1, 3, 9. Among these, $1\equiv 1,\; 3\equiv 3,\; 9\equiv 4\pmod{5}$. Only $n_5=1$ works. So there is exactly \textbf{1} Sylow 5-subgroup (and it is normal).

\item (a). For Sylow 5-subgroups: $n_5\mid 6$ and $n_5\equiv 1\pmod{5}$. Divisors of 6: 1, 2, 3, 6. Only $n_5=1$ or $n_5=6$. Since $n_5=6$ would account for $6\cdot 4=24$ elements of order 5, and combined with elements from other Sylow subgroups this is feasible, so $n_5=1$ or 6. However, by a more careful argument using the constraint that $G$ has 30 elements, one can show $n_5=1$. For $n_3$: $n_3\mid 10$ and $n_3\equiv 1\pmod{3}$, so $n_3\in\{1,10\}$; $n_3=10$ is possible. (c) is false; $S_3\times\mathbb{Z}/5\mathbb{Z}$ has order 30 and is not abelian. So only (a) is guaranteed.

\item The zero divisors in $\mathbb{Z}/12\mathbb{Z}$ are the non-zero elements that are not coprime to 12: $\bar{2}, \bar{3}, \bar{4}, \bar{6}, \bar{8}, \bar{9}, \bar{10}$. (E.g., $\bar{2}\cdot\bar{6}=\bar{0}$, $\bar{3}\cdot\bar{4}=\bar{0}$, $\bar{9}\cdot\bar{4}=\bar{0}$, $\bar{8}\cdot\bar{3}=\bar{0}$, $\bar{10}\cdot\bar{6}=\bar{0}$.)

\item (a), (c), and (d). $\mathbb{Z}$ is an integral domain. $\mathbb{Z}/6\mathbb{Z}$ has zero divisors ($\bar{2}\cdot\bar{3}=\bar{0}$). $\mathbb{Z}/7\mathbb{Z}$ is a field (hence integral domain). $\mathbb{Z}[x]$ is an integral domain since $\mathbb{Z}$ is.

\item $(x)$ is prime but not maximal. $\mathbb{Z}[x]/(x)\cong\mathbb{Z}$, which is an integral domain but not a field. So $(x)$ is prime (not maximal).

\item $\mathbb{Z}/(5)\cong\mathbb{Z}/5\mathbb{Z}$, which is the field with 5 elements $\{0,1,2,3,4\}$ with addition and multiplication mod~5.

\item No, $(2,x)$ is \textbf{not} principal. Suppose $(2,x)=(f(x))$. Then $f\mid 2$ and $f\mid x$ in $\mathbb{Z}[x]$. $f\mid 2$ means $f\in\{\pm 1,\pm 2\}$. If $f=\pm 1$ then $(f)=\mathbb{Z}[x]$, but $(2,x)\neq\mathbb{Z}[x]$ (e.g., $1\notin(2,x)$ since $2f(x)+xg(x)$ evaluated at $x=0$ gives an even number). If $f=\pm 2$ then $x\notin(2)$. Contradiction. So $(2,x)$ is not principal, and $\mathbb{Z}[x]$ is not a PID.

\item $\operatorname{char}(\mathbb{F}_9)=3$, since $|\mathbb{F}_9|=3^2$.

\item Yes. Over $\mathbb{F}_2$, a cubic is irreducible iff it has no root. Check: $f(0)=0+0+1=1\neq 0$, $f(1)=1+1+1=1\neq 0$. No roots, so $x^3+x+1$ is irreducible over $\mathbb{F}_2$.

\item $[\mathbb{Q}(\sqrt[3]{2}):\mathbb{Q}]=3$. The minimal polynomial of $\sqrt[3]{2}$ over $\mathbb{Q}$ is $x^3-2$, which is irreducible by Eisenstein with $p=2$.

\item $36 = 2^2\cdot 3^2$. By the Fundamental Theorem of Finite Abelian Groups, we decompose independently: for $2^2$: $\mathbb{Z}/4\mathbb{Z}$ or $(\mathbb{Z}/2\mathbb{Z})^2$ (2 choices); for $3^2$: $\mathbb{Z}/9\mathbb{Z}$ or $(\mathbb{Z}/3\mathbb{Z})^2$ (2 choices). Total: $2\times 2=\mathbf{4}$ abelian groups.

\end{enumerate}

\end{document}
